\begin{abstract}

地球系统因其极广的时空尺度、极高的状态空间维度以及混沌效应与随机扰动的深度耦合，基于传统理论难以深入理解其中发生的复杂过程。
本文发展了一套统一的数学框架，从微分几何视角，将地球系统描述为流形上的随机动力学过程。
我们提出物理约束将系统状态限制在远低于环境空间维度的光滑流形上 (流形假设)，为经验正交函数分解 (EOF)、自编码机 (AE) 等降维方法提供了解释。
地球系统状态的演化遵循流形上的Itô随机微分方程，其长期统计性质由Fokker-Planck方程导出的不变概率测度所刻画，随机轨迹和概率测度分别代表了天气和气候。
通过追踪初值不确定性的传播，我们从动力学方程导出了流形的度量张量，该度量自然编码了系统的可预报性极限: 沿Lyapunov不稳定方向，度量随预报时间线性增长，表征混沌效应；沿耗散方向，度量指数衰减至零，对应流形上的吸引子；沿准守恒模态，度量保持常数，支撑长期预报。
基于流形吸引子的拓扑结构，我们将气候变化分类为测度平滑演化 (I类) 与吸引子拓扑突变 (II类)，分别讨论了极端事件频率的指数增长与临界转变的预警信号。
另外，我们提出可以使用扩散模型进行数据驱动的流形学习: 扩散模型能够隐式地建模流形的几何结构。
概率流形理论可作为气候预测、临界预警、物理-AI融合建模的数学基础。

\end{abstract}
